\documentclass[11pt]{article}
\usepackage[a4paper, top=18mm, bottom=18mm, left=20mm, right=20mm]{geometry}
\usepackage[T2A]{fontenc}
\usepackage[utf8]{inputenc}
\usepackage[ukrainian]{babel}
\usepackage{graphicx}
\usepackage[export]{adjustbox}
\usepackage{amsmath}
\usepackage{amssymb}
\graphicspath{{/app/Quiz/Scripts/2023/ProgAll/15BinTree/}{/app/Quiz/Scripts/2021/Mod2021_3/BinTree/}{/app/Quiz/Scripts/2019/Mod2019_3/BinTree/}}
\setlength{\parindent}{0pt}
\setlength{\parskip}{4pt}
\pagestyle{empty}
\begin{document}
%begin
{\large\textbf{Контрольна робота}}

\textbf{Варіант \No\,1}\hfill \textit{ПІБ:}\,\underline{\hspace{60mm}}
\vspace{4mm}

%beginitem
1. 
try:
    res = 3**4-2
    if isinstance(res, bool):
        print(f'bool;{res}')
    elif isinstance(res, int):
        print(f'int;{res}')
    elif isinstance(res, float):
        print(f'float;{res:.2f}')
    else:
        print('???')
except:
    print('Z')

%enditem

%beginitem
2. 
Записати ВИРАЗ, значенням якого є множина, що складається з цілих чисел від 19 до 2011 (включно), що діляться на 7.
%enditem

%beginitem
3. Що буде виведено за виконання фрагменту коду?

%code
print('R', end=' ')
d=['0','1','2','3','4']
d[1:1]=[1,0]
print(d)
%code


%enditem

%beginitem
4. %goodpath:Scripts/2018/Mod2018_1/01-good.txt
Відмітити все, що є однією правильною лексемою:
( 1 ) <>

( 2 ) FIFO

( 3 ) >

( 4 ) =<

%enditem

%beginitem
5. 
Записати ВИРАЗ, значенням якого є список, що складається з квадратних коренівцілих чисел від 19 до 2011 (включно), що діляться на 7, записаних в обернено\-му порядку.

%enditem

%beginitem
6. 

print('R', end=' ')
for d in range(-4,-4,3):
    if d>=6:
        continue
        print(d, end=' ')
        d=1
    if d>=6:
        break
    else:
        print(d, end=' ')
    print(d)

%enditem

%beginitem
7. 
try:
    for d in range(-4, -4, 3):
        if d>=6:
            continue
        print(d, end=';');d = 1
        if d>=6:
            break
    else:
        print(d,end=';')
    print(d)
except:
    print('Z')

%enditem

%beginitem
8. 
Знайти кількість відношень на n-елементній множині, які неантирефлексивні або неантисиметричні
%enditem

%beginitem
9. 
Записати ВИРАЗ, значенням якого є множина, що складається з квадратних коренівцілих чисел від 19 до 2011 (включно), що діляться на 7.
%enditem

%beginitem
10. 
Записати ВИРАЗ, значенням якого є список, що складається з цілих чисел від 19 до 2011 (включно), причому спочатку за спаданням записано всі, що діляться на 7, а потім решту за спаданням.
%enditem

%beginitem
11. 
try:
    res = 7//7%7*7
    if isinstance(res, bool):
        print(f'bool;{res}')
    elif isinstance(res, int):
        print(f'int;{res}')
    elif isinstance(res, float):
        print(f'float;{res:.2f}')
    else:
        print('???')
except:
    print('ERR')

%enditem

%beginitem
12. 
Записати ВИРАЗ, значенням якого є список, що складається з цілих чисел від 19 до 2011 (включно), причому спочатку за спаданням записано всі, що діляться на 7, а потім решту за спаданням.

%enditem

%beginitem
13. 
try:
    print(8, end=';')
    print(int('4'), end=';')
    print(5, end=';')
except KeyboardInterrupt:
    print(7, end=';')
except Exception:
    print(9, end=';')
except:
    print('ERR', end=';')
else:
    print(6, end=';')
finally:
    print(1)

%enditem

%beginitem
14. Що буде виведено за виконання фрагменту коду?

%code
#include <iostream>
using namespace std;

bool f(int n){
    cout<<"f";
    return n>=0;
}

int main(){
    cout<<(f(2) || f(6));
    return 0;
}
%code

%enditem

%beginitem
15. Що буде виведено за виконання фрагменту коду?
%code
for d in range(-4,-4,3):
    if d>=6:
        continue
        print(d,end=' ')
        d=1
    if d>=6:
        break
else:
    print(d, end=' ')
print(d, end=' ')
%code

%enditem

%beginitem
16. 

print('R', end=' ')
d=['0','1','2','3','4']
d.append('13')
print(d)


%enditem

%beginitem
17. 

a,b,c=8,4,5
def g(b):
global c    a-=5
    b=1
    c=3
    return a+b+c

a,b,c=7,9,6
print(g(b),a,b,c)

%enditem

%beginitem
18. 

def g(a,b):
    c=59
    if b:
        c=8
    elif b>=4:
         return 4
    else: 
        return 5
    return c

print(g(2,2))
%enditem

%beginitem
19. %goodpath:Scripts/2019/Mod2019_1/Lex/03-good.txt
Відмітити все, що є коректним оператором С++:
( 1 ) n=2; n--

( 2 ) 3**-1

( 3 ) x=float("inf")

( 4 ) False=0

%enditem

%beginitem
20. 
try:
    n = 14
    while n>=8:
        n += -2
    print(n, end=';')
except:
    print('ERR')

%enditem

%beginitem
21. Що буде виведено за виконання фрагменту коду?
%code
#include <iostream>
using namespace std;

int f(int &x, int &y){
    x = 5;
    y-= 9;
    return y;
}

int main(){
    int a = 3, b = 2;
    b = f(b, b);
    cout << a << ":" << b <<':';
    {
        int a = 4, b = 7;
        cout << ((b>=7) || ((a-=2) >= 7))
             << a << ":" << b << ":";
    }
    {
        inta = 8;
        cout << a << ':';
    }
    cout << a << ":" << b << endl;
    return 0;
}
%code


%enditem

%beginitem
22. 
Записати ВИРАЗ, значенням якого є список, що складається з цілих чисел від 19 до 2011 (включно), причому спочатку за спаданням записано всі, що діляться на 7, а потім решту за спаданням.
%enditem

%beginitem
23. 
Змінні $next$ та $prev$ вузла списку позначають наступний та попередній за ним вузол відповідно. Починаючи з голови, у список записано послідовні цілі числа від $-21$ до $36$. Нехай змінна $p$ позначає вузол, що зберігає значення $2$.
Яке значення зберігається у вузлі $p.next.next.next.next.next$ ?

%enditem

%beginitem
24. 

\begin{center}
	\adjustimage{max size={0.8\linewidth}{0.8\paperheight}}
	{../15BinTree/trees_exp/155.png}
\end{center}
Указати послідовність міток вершин дерева, зображеного на рис., що утвориться в результаті обходу дерева в оберненому порядку. Мітки записувати через пробіл.\par Алгоритм починає роботу з кореня (на рис. це верхня вершина).
%answer:155 оберненому

%enditem

%beginitem
25. Що буде виведено за виконання фрагменту коду?

%code
print('R', end=' ')
d = ('a','b','c','d','e',0,1,2,3,4)
print(d[-2])
%code


%enditem

%beginitem
26. 
Змінні next та prev вузла списку позначають наступний та попередній за ним вузол відповідно. Починаючи з голови у список записано послідовні цілі числа від -21 до 36. Нехай змінна p позначає вузол, що зберігає значення 2.
Яке значення зберігається у вузлі p.next.next.next.next.next ?
%enditem

%beginitem
27. Що буде виведено за виконання фрагменту коду?

%code
print('R', end=' ')
d=['0','1','2','3','4']
d.append('13')
print(d)
%code


%enditem

%beginitem
28. Що буде виведено за виконання фрагменту коду?

%code
#include <iostream>
using namespace std;

int g(int a, int b){
    int c = 59;
    if (b)
        c = 8;
    else if (b >= 4)
         return 4;
    else 
        return 5;
    return c;
}

int main(){
    cout << g(2, 2);
    return 0;
}
%code

%enditem

%beginitem
29. 
for d in range(-4,-4,3):
     if d>=6: continue
     print(d,end=' ');d=1
     if d>=6: break
   else: print(d,end=' ')
   print(d)
%enditem

%beginitem
30. 

x,y,z=3,4,5
x,y,z,z=y,x,9,6
print(x,y,z)

%enditem

%beginitem
31. Що буде виведено за виконання фрагменту коду?

%code
print("5j">="5.0j")
%code

%enditem

%beginitem
32. 
try:
    res = not5.0>=5 or notFalse==6 or 7!=6.0
    if isinstance(res, bool):
        print(f'bool;{res}')
    elif isinstance(res, int):
        print(f'int;{res}')
    elif isinstance(res, float):
        print(f'float;{res:.2f}')
    else:
        print('???')
except:
    print('ERR')

%enditem

%beginitem
33. Що буде виведено за виконання фрагменту коду?

%code
#include <iostream>
using namespace std;

int a = 8, b = 4, c = 5;

int g(int &b){
int c;    a = 5;
    b = 1;
    c = 3;
    return a + b + c;
}

int main(){
    inta = 7;
    intb = 9;
    intc = 6;
    cout << g(b) << ':';
    cout << a << ':' << b << ':' << c;
    return 0; 
}
%code

%enditem

%beginitem
34. 

def g(d):
    x=59
    if d: 
        x=8
    elif d>=4:
         return 4
    else:
         return 5
    return x

print(g(2))

%enditem

%beginitem
35. 
Скільки 2017-елементних підмножин множини натуральних від 1 до 20172017 містять рівно два числа, що менше 172010
%enditem

%beginitem
36. 
7//7%7*7
%enditem

%beginitem
37. Що буде виведено за виконання фрагменту коду?

%code
print('R', end=' ')
d=['0','1','2','3','4','5','6','7','8','9']
print(d[-4::3])
%code


%enditem

%beginitem
38. 

\begin{center}
	\adjustimage{max size={0.8\linewidth}{0.8\paperheight}}
	{../BinTree/trees_exp/155.png}
\end{center}
Указати послідовність міток вершин дерева, зображеного на рис., що утвориться в результаті обходу дерева в оберненому порядку. Мітки записувати через пробіл.\par Алгоритм починає роботу з кореня (на рис. це верхня вершина).
%answer:155 оберненому

%enditem

%beginitem
39. Що буде виведено за виконання фрагменту коду?

%code
def g(a,b):
    c=59
    if b:
        c=8
    elif b>=4:
         return 4
    else: 
        return 5
    return c

print(g(2,2))
%code

%enditem

%beginitem
40. Що буде виведено за виконання фрагменту коду?

%code
print('R', end=' ')
d = ['0','1','2','3','4','5','6','7','8','9']
print(d[-4::3])
%code


%enditem

%beginitem
41. Що буде виведено за виконання фрагменту коду?
%code

a,b,c=9,6,7
def g(a,b=8,c=6):
    print(a,b,c,end="")

g(3,4,b=5)
print(a,b,c)

%code

%enditem

%beginitem
42. 
Записати ВИРАЗ, значенням якого є множина, що складається з квадратних коренівцілих чисел від 19 до 2011 (включно), що діляться на 7.
%enditem

%beginitem
43. 
Записати ВИРАЗ, значенням якого є словник, ключами якого є цілі числа від 19 до 2011 (включно), що діляться на 7, а ключам відповідають їх куби 
%enditem

%beginitem
44. 
n=14
   while n>=8: n+=-2
   print(n)
%enditem

%beginitem
45. Що буде виведено за виконання фрагменту коду?

%code
n=14
while n>=8:
    n+=-2
print(n)
%code


%enditem

%beginitem
46. 
char d[]="abcde0124"; cout<<sizeof(d)<<','<<d<<endl;
   d[4]='\0'; cout<<sizeof(d)<<','<<d<<'*';
%enditem

%beginitem
47. 
Рядки журналу через = містять ім'я користувача (ідентифікатор), час події,тип події, код події, наприклад
\begin{verbatim}
s="username = 2022-05-05 13:26 = ERROR = 404"
\end{verbatim}
у коді 
\begin{verbatim}
res = re.fullmatch('???', s) 
s = ''
code_str = ??? 
\end{verbatim}
чим треба замінити кожен з ???, щоб значенням code\_str був код події? 



%enditem

%beginitem
48. Що буде виведено за виконання фрагменту коду?

%code
#include <iostream>

bool f(int n){
    std::cout<<"f";
    return n>=0;
}

int main(){
    std::cout<<(f(2) || f(6));
    return 0;
}
%code

%enditem

%beginitem
49. 
try:
    n = 8
    while n<=9:
        n += 2
    print(n, end=';')
except:
    print('ERR')

%enditem

%beginitem
50. Що буде виведено за виконання фрагменту коду?

%code
cout << (5.0 >= 5 == false != 6);
%code

%enditem

%beginitem
51. 
Записати ВИРАЗ, значенням якого є множина, що складається з цілих чисел від 19 до 2011 (включно), що діляться на 7.

%enditem

%beginitem
52. 
a = 8
b = 4
c = 5

def g(b):
global c    a -= 5
    b = 1
    c = 3
    return a + b + c

a = 7
b = 9
c = 6

try:
    print(g(b), a, b, c, sep=';')
except:
    print('Z', a, b, c, sep=';')

%enditem

%beginitem
53. Що буде виведено за виконання фрагменту коду?

%code
print(5.0 == 5 != False > 6)
%code

%enditem

%beginitem
54. %goodpath:Scripts/2018/Mod2018_1/03-good.txt
Відмітити все, що є коректним оператором С++:
( 1 ) int n,n;

( 2 ) int f(int n);

( 3 ) int x='a'-'c';

( 4 ) x=*2;

%enditem

%beginitem
55. 
Записати ВИРАЗ, значенням якого є множина, що складається з цілих чисел від 19 до 2011 (включно), що діляться на 7.
%enditem

%beginitem
56. 
Рядки журналу через двокрапку містять ім'я користувача (ідентифікатор), час події,тип події, код події, наприклад
\begin{verbatim}
s="username : 2022-05-05 13:26 : ERROR : 404"
\end{verbatim}
у коді 
\begin{verbatim}
res = re.fullmatch('???', s) 
s = ''
code_str = ??? 
\end{verbatim}
чим треба замінити кожен з ???, щоб значенням code\_str був код події? 



%enditem

%beginitem
57. 

\begin{center}
	\adjustimage{max size={0.8\linewidth}{0.8\paperheight}}
	{../BinTree/trees_exp/49.png}
\end{center}
Указати послідовність міток вершин дерева, зображеного на рис., що утвориться в результаті обходу дерева в оберненому порядку. ̳тки записувати через пробіл.\par Обхід дерева починається з кореня (на рис. це верхня вершина).
%answer:49 оберненому
%enditem

%beginitem
58. 
"5j">="5.0j"
%enditem

%beginitem
59. 
def f(n):
    print('f', end='')
    return n>=0

try:
    print(f(2) or f(6))
except:
    print('ERR')

%enditem

%beginitem
60. Що буде виведено за виконання фрагменту коду?

%code
print(3**4-2)
%code

%enditem

%beginitem
61. Що буде виведено за виконання фрагменту коду?

%code
print(not5.0>=5 or notFalse==6 or 7!=6.0)
%code

%enditem

%beginitem
62. 
for d in range(-4,-4,3):
    if d>=6:
        continue
        print(d,end=' ')
        d=1
    if d>=6:
        break
else:
    print(d, end=' ')
print(d, end=' ')

%enditem

%beginitem
63. 
d=['0','1','2','3','4']; d.append('13'); print(d)
%enditem

%beginitem
64. 
У папці X31 розташовано проект, до якого входить синаксично правильна одиниця трансляції 1.cpp з рядком\par
{\#}include "2.h"
\parФайла 2.h у папці X31 нема, але він є в папці Y7. Що треба зробити, щоб одиниця трансляції 1.cpp, могла успішно відкомпілюватися?Переміщувати, копіювати та змінювати файли з кодом заборонено.\par\vspace{1.5cm}
%enditem

%beginitem
65. Що буде виведено за виконання фрагменту коду?

%code
cout << (14 / 14 / 14 / 14);
%code

%enditem

%beginitem
66. 

try:
    print(8,end="")
    print(int("4"),end="")
    print(5,end="")
except TypeError: print(7,end="")
except ZeroDivisionError: print(9,end="")
else: print(6,end="")
finally: print(1,end="")

%enditem

%beginitem
67. 
Задано тип вузла списку
struct Node \{int n; Node *prev, *next;\};
У списку зберігаються послідовні цілі числа від -21 до 36. Вказівник p встановлено на вузол, в якому зберігається число 2.
Чому дорівнює значення p->next->next->next->next->next->n ?
%enditem

%beginitem
68. Що буде виведено за виконання фрагменту коду?

%code
d = ('a','b','c','d','e',0,1,2,3,4)
print(d[-2])
%code


%enditem

%beginitem
69. Що буде виведено за виконання фрагменту коду?

%code
n=8
while n<=9:
    n+=2
print(n, end=' ')
%code


%enditem

%beginitem
70. %goodpath:Scripts/2019/Mod2019_1/Lex/01-good.txt
Відмітити все, що є однією правильною лексемою:
( 1 ) x*y

( 2 ) {

( 3 ) "a\"a"

( 4 ) =>

%enditem

%beginitem
71. Що буде виведено за виконання фрагменту коду?

%code
"5j">="5.0j"
%code

%enditem

%beginitem
72. Що буде виведено за виконання фрагменту коду?
%code
if (10 >= 10)
    cout << "v";
else
    cout << "v";
%code


%enditem

%beginitem
73. Що буде виведено за виконання фрагменту коду?

%code
print('R', end=' ')
d=('a','b','c','d','e',0,1,2,3,4)
print(d[-2])
%code


%enditem

%beginitem
74. Що буде виведено за виконання фрагменту коду?

%code
print(7//7%7*7)
%code

%enditem

%beginitem
75. 
d=('a','b','c','d','e',0,1,2,3,4); print(d[-2])
%enditem

%beginitem
76. Що буде виведено за виконання фрагменту коду?

%code
def g(d):
    x=59
    if d: 
        x=8
    elif d>=4:
         return 4
    else:
         return 5
    return x

print(g(2))
%code

%enditem

%beginitem
77. 
Записати ВИРАЗ, значенням якого є список, що складається з квадратних коренівцілих чисел від 19 до 2011 (включно), що діляться на 7, записаних в оберненому порядку.
%enditem

%beginitem
78. 
code
d = 0
m = 1

class D:
    m = 2
    
    def __init__(self):
global m        self.m = 3
        m = 4
        D.m = 5

obj = D()
print(d, m, D.m, obj.m, sep=' ')
code

%enditem

%beginitem
79. Що буде виведено за виконання фрагменту коду?

%code
cout << 14 / 14 / 14 / 14;
%code

%enditem

%beginitem
80. 
def f(arg, n):
    n = 8
    tmp = arg
    tmp[1:1] = [1,0]
    return len(tmp)>3

try:
    d = ['0','1','2','3','4']
    n = 4
    f(d, n)
    print(n, end=';')
    for el in d:
        print(el, end=';')
except:
    print('ERR')

%enditem

%beginitem
81. Що буде виведено за виконання фрагменту коду?

%code
a = 8
b = 4
c = 5

def g(b):
global c    a=5
    b=1
    c=3
    return a+b+c

a, b, c = 7,9,6
try:
    print(g(b), a, b, c, sep=';')
except:
    print('Z', a, b, c, sep=';')
%code

%enditem

%beginitem
82. Що буде виведено за виконання фрагменту коду?

%code
a = 8
b = 4
c = 5

def g(b):
global c    a-=5
    b=1
    c=3
    return a+b+c

a, b, c = 7, 9, 6
try:
    print(g(b), a, b, c, sep=';')
except:
    print('Z', a, b, c, sep=';')
%code

%enditem

%beginitem
83. Що буде виведено за виконання фрагменту коду?

%code
try:
    n = 8
    while n<=9:
        n += 2
    print(n, end=';')
except:
    print('Z')

%code

%enditem

%beginitem
84. 
Проект складається з од���ї одиниці трансляції 1.cpp, яка починається таким фрагментом\par
long g(){ return myFunc();}\par
Функція myFunc(); означена в 2.cpp.\par
Що треба зробити для того, щоб 1.cpp успішно відкомпілювався?\par\vspace{1.5cm}
Що треба зробити для того, щоб за проектом зібрався виконуваний код?\par

%enditem

%beginitem
85. 

def f(n):
    print("f", end="");
    return n>=0

print(f(2) or f(6))

%enditem

%beginitem
86. Що буде виведено за виконання фрагменту коду?
%Оригинал был утерян. Требует отладки!
%code
__d = 0

class D:
    __d = 1
    
    def __init__(self):
        self.__d = 2
        __d = 3
        D.__d = 4

obj = D()
D.__d = 5
try:
    print(__d, end=' ')
    print(obj.__d, end=' ')
    print(D.__d, end=' ')
    print(obj.__d, end=' ')
    print(obj._D__d)
except:
    print('ERR')
%code

%enditem

%beginitem
87. Що буде виведено за виконання фрагменту коду?

%code
print('R', end=' ')
n=14
while n>=8:
    n+=-2
print(n)
%code


%enditem

%beginitem
88. 
Кожен рядок текстового файлу містить по два дійсних числа, розділені білими символами.
Написати функцію, що для заданого текстового файлу будує новий текстовий файл, рядки якого крім зображень дйсних чисел ще містять результат обчислення на них функції $f$. У вихідному файлі зображення чисел мають розділятися пробілами. 
$$f(x, y) =\frac{\pi x - 26}{(\arcsin x - 0.5) (y - 22)}$$ 
Якість коду також оцінюється.


%enditem

%beginitem
89. Що буде виведено за виконання фрагменту коду?

%code
a,b,c=8,4,5
def g(b):
global c    a-=5
    b=1
    c=3
    return a+b+c

a,b,c=7,9,6
print(g(b),a,b,c)
%code

%enditem

%beginitem
90. 

print('R', end=' ')
d=['0','1','2','3','4','5','6','7','8','9']
print(d[-4::3])


%enditem

%beginitem
91. 
a,b,c=8,4,5
    def g(b):
      global c      a=5
      b=1
      c=3
      return a+b+c
    a,b,c=7,9,6
    print(g(b),a,b,c)
%enditem

%beginitem
92. 
n=8
   while n<=9: n+=2
   print(n)
%enditem

%beginitem
93. Що буде виведено за виконання фрагменту коду?
%code
#include <iostream>
using namespace std;

int f(int &x, int &y){
    x = 5;
    y-= 9;
    return y;
}

int main(){
    int a = 3, b = 2;
    b = f(b, b);
    cout << a << ":" << b <<':';
    {
        int a = 4, b = 7;
        cout << ((b>=7) || ((a-=2) >= 7)) << ':';
        cout << a << ":" << b << ':';
    }
    {
        inta = 8;
        cout << a << ':';
    }
    cout << a << ":" << b << endl;
    return 0;
}
%code


%enditem

%beginitem
94. 
Записати ВИРАЗ, значенням якого є словник, ключами якого є цілі числа від 19 до 2011 (включно), що діляться на 7, а ключам відповідають їх куби 

%enditem

%beginitem
95. 

print('R', end=' ')
d=('a','b','c','d','e',0,1,2,3,4)
print(d[-2])


%enditem

%beginitem
96. Що буде виведено за виконання фрагменту коду?
%code
if (10 >= 10)
    cout << "v";
else
    cout << "o";
%code


%enditem

%beginitem
97. 
a = 8
b = 4
c = 5

def g(b):
global c    a = 5
    b = 1
    c = 3
    return a + b + c

a = 7
b = 9
c = 6

try:
    print(g(b), a, b, c, sep=';')
except:
    print('ERR', a, b, c, sep=';')

%enditem

%beginitem
98. Що буде виведено за виконання фрагменту коду?

%code
d=['0','1','2','3']
d.append('13')
print(d)
%code


%enditem

%beginitem
99. Що буде виведено за виконання фрагменту коду?
%code

def f(n):
    print("f", end="");
    return n>=0

print(f(2) or f(6))
%code

%enditem

%beginitem
100. Що буде виведено за виконання фрагменту коду?

%code
a,b,c=8,4,5
def g(b):
    a=5
    b=1
    c=3
    return a+b+c

a,b,c=7,9,6
print(g(b),a,b,c)
%code

%enditem

%end
\newpage
\end{document}

